\documentclass[11pt,a4paper]{article}

% Idioma e fontes (compilar com XeLaTeX)
\usepackage{fontspec}
\usepackage[brazilian]{babel}

\usepackage[margin=2.5cm]{geometry}
\usepackage{amsmath, amssymb}
\usepackage{graphicx}
\usepackage{enumitem}
\usepackage[hidelinks]{hyperref}

% Dados da tarefa
\newcommand{\curso}{Nome do curso}
\newcommand{\tarefa}{Tarefa 1}
\newcommand{\estudante}{Nome Sobrenome}
\newcommand{\professor}{Nome do professor}

% Questão numerada e ambiente de resposta
\newcounter{questao}
\newcommand{\questao}[1]{%
  \stepcounter{questao}%
  \section*{Questão \thequestao\ifx\relax#1\relax\else: #1\fi}}
\newenvironment{resposta}{\par\noindent\textbf{Resposta.}\ }{\hfill$\blacksquare$\par}

\begin{document}

\noindent
\begin{minipage}[t]{0.6\textwidth}
  \textbf{\curso}\\
  \tarefa
\end{minipage}%
\begin{minipage}[t]{0.4\textwidth}
  \raggedleft
  \estudante\\
  Professor: \professor\\
  \today
\end{minipage}
\par\medskip\hrule\medskip

\questao{Derivadas}
Calcule $\dfrac{d}{dx}\left(x^2 \sin x\right)$.

\begin{resposta}
  Pela regra do produto,
  \[
    \frac{d}{dx}\left(x^2 \sin x\right) = 2x \sin x + x^2 \cos x.
  \]
\end{resposta}

\questao{}
Enuncie duas propriedades da integral definida.

\begin{resposta}
  \begin{enumerate}[label=(\alph*)]
    \item Linearidade: $\int_a^b (\alpha f + \beta g) = \alpha \int_a^b f + \beta \int_a^b g$.
    \item Aditividade: $\int_a^b f = \int_a^c f + \int_c^b f$.
  \end{enumerate}
\end{resposta}

\end{document}
